\documentclass[12pt]{article}
\usepackage{amsmath, amssymb}
\usepackage{geometry}
\geometry{margin=1in}
\setlength{\parskip}{0.6em}
\setlength{\parindent}{0pt}

\begin{document}

\section*{Solution Key: Human Capital Investment Decision}

\subsection*{Given Information}

Johny is considering a \textbf{3-year college diploma}.

\begin{itemize}
    \item Length of program: \(3\) years
    \item Annual tuition: \(\$1{,}200\)
    \item Annual books/materials: \(\$300\)
    \item Earnings without college: \(\$8{,}000\) per year
    \item Earnings with diploma: \(\$14{,}000\) per year
    \item Working life after graduation: \(12\) years
    \item Market interest rate: \(r = 6\% = 0.06\)
\end{itemize}

\hrule
\vspace{0.5em}

\subsection*{1. Total annual direct cost of attending college}

The annual \textbf{direct cost} is:

\[
\text{Direct Cost} = \text{Tuition} + \text{Books}
\]

\[
\text{Direct Cost} = 1200 + 300 = 1500
\]

\[
\boxed{\text{Annual direct cost} = \$1{,}500}
\]

\hrule
\vspace{0.5em}

\subsection*{2. Opportunity cost per year of attending college}

The \textbf{opportunity cost} is the income Johny gives up by going to school instead of working.

\[
\text{Opportunity Cost per year} = \$8{,}000
\]

\[
\boxed{\text{Opportunity cost per year} = \$8{,}000}
\]

\hrule
\vspace{0.5em}

\subsection*{3. Total cost of the investment over the 3-year period (ignoring discounting)}

The total annual cost is:

\[
\text{Annual Total Cost} = \text{Direct Cost} + \text{Opportunity Cost}
\]

\[
\text{Annual Total Cost} = 1500 + 8000 = 9500
\]

Since the program lasts 3 years:

\[
\text{Total Cost over 3 years} = 3 \times 9500 = 28{,}500
\]

\[
\boxed{\text{Total undiscounted cost over 3 years} = \$28{,}500}
\]

\hrule
\vspace{0.5em}

\subsection*{4. Annual earnings gain from completing the diploma}

The additional annual earnings from the diploma are:

\[
\Delta Y = 14{,}000 - 8{,}000 = 6{,}000
\]

\[
\boxed{\text{Annual earnings gain} = \$6{,}000}
\]

\hrule
\vspace{0.5em}

\subsection*{5. Present value of benefits}

Johny receives the additional earnings of \(\$6{,}000\) each year for \(12\) years after graduation.

So the present value of benefits, measured \textbf{at the time of graduation}, is the present value of a 12-year annuity:

\[
PV_B^{(t=3)} = 6000 \left( \frac{1-(1.06)^{-12}}{0.06} \right)
\]

First compute:

\[
(1.06)^{12} \approx 2.0122
\]

\[
(1.06)^{-12} \approx \frac{1}{2.0122} \approx 0.49697
\]

Then:

\[
\frac{1-0.49697}{0.06} = \frac{0.50303}{0.06} \approx 8.3838
\]

So:

\[
PV_B^{(t=3)} = 6000(8.3838) \approx 50{,}302.8
\]

This is the present value \textbf{at the time Johny finishes school}. Since the schooling decision is made \textbf{today}, we must discount this amount back 3 years:

\[
PV_B^{(t=0)} = \frac{50{,}302.8}{(1.06)^3}
\]

\[
(1.06)^3 = 1.191016
\]

\[
PV_B^{(t=0)} \approx \frac{50{,}302.8}{1.191016} \approx 42{,}235.9
\]

\[
\boxed{PV(\text{Benefits}) \approx \$42{,}236}
\]

\hrule
\vspace{0.5em}

\subsection*{6. Present value of costs}

Johny incurs total annual cost of \(\$9{,}500\) for 3 years. These costs occur during years 1, 2, and 3.

Thus the present value of costs is:

\[
PV_C = \frac{9500}{1.06} + \frac{9500}{(1.06)^2} + \frac{9500}{(1.06)^3}
\]

Compute each term:

\[
\frac{9500}{1.06} \approx 8962.26
\]

\[
\frac{9500}{1.1236} \approx 8454.97
\]

\[
\frac{9500}{1.191016} \approx 7974.71
\]

Add them:

\[
PV_C \approx 8962.26 + 8454.97 + 7974.71 = 25{,}391.94
\]

\[
\boxed{PV(\text{Costs}) \approx \$25{,}392}
\]

\hrule
\vspace{0.5em}

\subsection*{7. Net Present Value (NPV) decision}

The net present value is:

\[
NPV = PV(\text{Benefits}) - PV(\text{Costs})
\]

\[
NPV = 42{,}235.9 - 25{,}391.94 = 16{,}843.96
\]

\[
\boxed{NPV \approx \$16{,}844}
\]

Since:

\[
NPV > 0
\]

Johny \textbf{should invest} in the diploma.

\[
\boxed{\text{Yes, Johny should undertake the investment.}}
\]

\hrule
\vspace{0.5em}

\subsection*{8. Effect of a higher interest rate (\(r=8\%\))}

If the market interest rate rises, future earnings are discounted more heavily.

This means:

\begin{itemize}
    \item The present value of future benefits falls.
    \item The attractiveness of education decreases.
\end{itemize}

Now compute the present value of benefits at \(8\%\):

\[
PV_B^{(t=3)} = 6000 \left( \frac{1-(1.08)^{-12}}{0.08} \right)
\]

\[
(1.08)^{12} \approx 2.5182
\]

\[
(1.08)^{-12} \approx 0.3971
\]

\[
\frac{1-0.3971}{0.08} = \frac{0.6029}{0.08} \approx 7.5363
\]

\[
PV_B^{(t=3)} \approx 6000(7.5363) = 45{,}217.8
\]

Discounting back 3 years:

\[
PV_B^{(t=0)} = \frac{45{,}217.8}{(1.08)^3}
\]

\[
(1.08)^3 = 1.2597
\]

\[
PV_B^{(t=0)} \approx \frac{45{,}217.8}{1.2597} \approx 35{,}895.7
\]

Now compute the present value of costs:

\[
PV_C = \frac{9500}{1.08} + \frac{9500}{(1.08)^2} + \frac{9500}{(1.08)^3}
\]

\[
\frac{9500}{1.08} \approx 8796.30
\]

\[
\frac{9500}{1.1664} \approx 8144.15
\]

\[
\frac{9500}{1.2597} \approx 7541.48
\]

\[
PV_C \approx 8796.30 + 8144.15 + 7541.48 = 24{,}481.93
\]

Thus:

\[
NPV_{8\%} \approx 35{,}895.7 - 24{,}481.93 = 11{,}413.77
\]

So the NPV is still positive, but smaller.

\[
\boxed{\text{A higher interest rate reduces the attractiveness of education, but Johny should still invest.}}
\]

\hrule
\vspace{0.5em}

\subsection*{9. Approximate internal rate of return (IRR)}

The IRR is the discount rate that makes:

\[
PV(\text{Benefits}) = PV(\text{Costs})
\]

We already know:

\begin{itemize}
    \item At \(6\%\), NPV is positive.
    \item At \(8\%\), NPV is still positive.
\end{itemize}

So the IRR must be \textbf{above 8\%}.

Try \(12\%\):

\[
PV_B^{(t=3)} = 6000\left(\frac{1-(1.12)^{-12}}{0.12}\right)
\]

\[
(1.12)^{12} \approx 3.89598
\]

\[
(1.12)^{-12} \approx 0.2567
\]

\[
\frac{1-0.2567}{0.12} = \frac{0.7433}{0.12} \approx 6.194
\]

\[
PV_B^{(t=3)} \approx 6000(6.194)=37{,}164
\]

Discount back 3 years:

\[
PV_B^{(t=0)} = \frac{37{,}164}{(1.12)^3}
\]

\[
(1.12)^3 = 1.404928
\]

\[
PV_B^{(t=0)} \approx 26{,}453
\]

Now costs at 12\%:

\[
PV_C = \frac{9500}{1.12} + \frac{9500}{(1.12)^2} + \frac{9500}{(1.12)^3}
\]

\[
\approx 8482.14 + 7573.34 + 6761.91 = 22{,}817.39
\]

So:

\[
NPV_{12\%} \approx 26{,}453 - 22{,}817 = 3{,}636
\]

Still positive, so IRR is above \(12\%\).

Try \(15\%\):

Benefits:
\[
PV_B^{(t=3)} = 6000\left(\frac{1-(1.15)^{-12}}{0.15}\right)
\]

\[
(1.15)^{12}\approx 5.3503,\qquad (1.15)^{-12}\approx 0.1869
\]

\[
\frac{1-0.1869}{0.15} \approx 5.4207
\]

\[
PV_B^{(t=3)} \approx 6000(5.4207)=32{,}524.2
\]

Discount back 3 years:

\[
(1.15)^3=1.520875
\]

\[
PV_B^{(t=0)}\approx \frac{32{,}524.2}{1.520875}\approx 21{,}385
\]

Costs:
\[
PV_C=\frac{9500}{1.15}+\frac{9500}{(1.15)^2}+\frac{9500}{(1.15)^3}
\]

\[
\approx 8260.87+7183.37+6246.41=21{,}690.65
\]

Now NPV is slightly negative:

\[
NPV_{15\%}\approx 21{,}385-21{,}691\approx -306
\]

So the IRR is between \(14\%\) and \(15\%\), very close to \(15\%\).

\[
\boxed{IRR \approx 14.8\% \text{ to } 15\%}
\]

\hrule
\vspace{0.5em}

\subsection*{10. Comparative statics}

\subsubsection*{(a) If expected post-education earnings fall}

Then:

\[
\Delta Y = \text{earnings with diploma} - \text{earnings without diploma}
\]

becomes smaller.

Therefore:
\begin{itemize}
    \item present value of benefits falls,
    \item NPV falls,
    \item IRR falls.
\end{itemize}

So Johny becomes \textbf{less likely} to invest.

\[
\boxed{\text{Lower post-education earnings reduce the attractiveness of schooling.}}
\]

\subsubsection*{(b) If tuition fees increase}

Higher tuition raises direct cost \(D\), so:
\begin{itemize}
    \item marginal cost rises,
    \item present value of costs rises,
    \item NPV falls.
\end{itemize}

Thus Johny becomes \textbf{less likely} to invest.

\[
\boxed{\text{Higher tuition reduces the incentive to invest in education.}}
\]

\subsubsection*{(c) If his working life becomes shorter}

A shorter working life means fewer years over which to earn the higher post-education income.

Thus:
\begin{itemize}
    \item present value of benefits falls,
    \item NPV falls,
    \item IRR falls.
\end{itemize}

So Johny becomes \textbf{less likely} to invest.

\[
\boxed{\text{A shorter working life reduces the return to education.}}
\]

\hrule
\vspace{0.8em}

\subsection*{Final Conclusion}

Johny should invest in the 3-year diploma because:

\[
PV(\text{Benefits}) > PV(\text{Costs})
\]

and

\[
NPV > 0
\]

The estimated internal rate of return is approximately:

\[
\boxed{IRR \approx 15\%}
\]

which is well above the market interest rate of \(6\%\).

\[
\boxed{\text{Therefore, the diploma is a worthwhile human capital investment.}}
\]

\end{document}